\chapter*{Introducción}\label{chapter:introduction}
\addcontentsline{toc}{chapter}{Introducción}

El \textit{Streptococcus pneumoniae}, conocido como neumococo, es la bacteria responsable de la enfermedad neumocócica. Este patógeno provoca un amplio espectro de infecciones, que van desde formas graves como la neumonía, la meningitis y la bacteriemia febril, hasta manifestaciones más frecuentes pero generalmente menos severas, como la otitis media, la sinusitis y la bronquitis \cite{cdc_pneumococica_2026, who_pneumococcal_2026, cdc_pneumococica_2024}.

Desde el punto de vista epidemiológico, la enfermedad neumocócica representa una elevada carga de mortalidad a nivel global, afectando especialmente a las edades extremas de la vida: niños menores de 5 años y adultos mayores de 65 años; con cientos de miles de muertes cada año \cite{who_pneumococcal_2026}. 

Dicho impacto ha impulsado un largo recorrido de investigación científica que se remonta a 1881, cuando Louis Pasteur y George Sternberg identificaron de forma independiente al neumococo. El estudio de esta bacteria contribuyó al desarrollo de la biología molecular, pues fue empleada, por ejemplo, en los experimentos de Avery, MacLeod y McCarty que demostraron el papel del ADN como material genético \cite{Neumoexpertos2016}. El interés por su control mediante vacunación fue creciendo a medida que se evidenciaba que la mortalidad persistía aun con tratamiento antibiótico disponible, lo que llevó a la autorización de la primera vacuna polisacárida en 1977 y de la primera vacuna conjugada en el año 2000 \cite{cdc_pneumococica_2024}.

Actualmente se conoce que el neumococo presenta una notable diversidad antigénica, con más de 90 serotipos identificados; sin embargo, tan solo 6 a 11 de ellos concentran aproximadamente el 70\% de las infecciones invasoras en niños. Esta distribución varía según la región geográfica, el grupo etario y el período analizado, lo que condiciona directamente la composición y la efectividad de las vacunas disponibles \cite{who_pneumococcal_2026, paho_neumococo_2026}.

Cuba ha desarrollado su propio programa de investigación y producción de vacunas antineumocócicas con el objetivo de hacer frente a la carga de la enfermedad en la población. En septiembre de 2024 se inició la aplicación de la vacuna conjugada de 7 valencias Quimi-Vio\textsuperscript{\textregistered}, dirigida a niños de uno a cinco años, y hasta la fecha se han inmunizado más de 95 mil niños en todo el país, con un muy buen perfil de seguridad \cite{prensa2025neumo11}. En ese mismo año también se incorporó al Programa de Inmunización Nacional la vacuna Pneumosil 10 valente, protegiendo a más de 27 500 lactantes contra las enfermedades derivadas del neumococo, con el respaldo financiero de la Alianza Mundial para las Vacunas y la Inmunización (GAVI) y el apoyo técnico de la Organización Panamericana de la Salud (OPS) y la Organización Mundial de la Salud (OMS) \cite{dpsalud_pneumosil_2024}. 

Más recientemente, se han iniciado ensayos clínicos de un nuevo candidato vacunal de 11 valencias: Quimi-Vio\textsuperscript{\textregistered} 11, destinado a adultos entre los 50 y 74 años de edad, en aras de ampliar la cobertura de protección a este grupo vulnerable \cite{prensa2025neumo11}, (véase anexos \ref{fig:bar_deaths}, \ref{fig:bar_sex} y \ref{fig:pie}). 

Evaluar la eficacia de Quimi-Vio 11 mediante un ensayo clínico convencional resulta extremadamente costoso, debido al elevado número de participantes que se requeriría. Sin embargo, es posible inferir dicha eficacia de forma más eficiente si se dispone de una vacuna similar ya aprobada y con datos de eficacia establecidos \cite{cruz2025}. 

A partir de este contexto, se formula la siguiente pregunta científica: ¿Cuál es la eficacia esperada del candidato vacunal Quimi-Vio\textsuperscript{\textregistered} 11 en adultos a partir de la evidencia de eficacia de Quimi-Vio\textsuperscript{\textregistered} 7 en niños? 

Para ello, se realizará un estudio de inmunopuente entre ambos grupos etarios, un enfoque estadístico y matemático que permite extrapolar la eficacia de una vacuna desde una población (niños) a otra (adultos) utilizando un marcador inmunológico capaz de estimar la protección inducida. 

Sobre esta base, el objetivo general de este trabajo es predecir la eficacia de Quimi-Vio\textsuperscript{\textregistered} 11 en población adulta, a partir de la eficacia observada en la población pediátrica para la vacuna Quimi-Vio\textsuperscript{\textregistered} 7.

Para alcanzar este objetivo, se plantean los siguientes objetivos específicos:
\begin{itemize} 
	\item Evaluar, mediante un enfoque de inmunopuente, la no inferioridad del candidato vacunal en adultos respecto a los serotipos en común con Quimi-Vio\textsuperscript{\textregistered} 7 en niños.
	\item Comparar métodos de estimación de parámetros de inmunogenicidad ante la presencia de un alto porcentaje de valores por debajo de los límites de cuantificación del ensayo clínico.
\end{itemize}

El presente trabajo se estructura en tres capítulos que permiten su comprensión y desarrollo de forma ordenada. En primer lugar, se abordan los fundamentos teóricos relacionados con la inmunología, se describe a detalle el concepto de puente inmunológico y sus antecedentes; así como los conceptos estadísticos fundamentales para realizar el estudio.

Posteriormente, se describen los métodos y procedimientos utilizados en la investigación; lo cual incluye el proceso de obtención y manejo de los datos, el cálculo de los diferentes parámetros de inmunogenicidad requeridos y el contraste de no inferioridad adoptado.

Finalmente, se presentan y discuten los resultados obtenidos, junto con sus principales limitaciones. Se incluyen, además, las referencias bibliográficas y los anexos con las figuras, tablas y gráficos obtenidos.
